\section{Institutional Background}\label{sec:background}

\subsection{Public procurement in S\~{a}o Paulo}

S\~{a}o Paulo's state government purchases common goods and services through the \emph{Bolsa Eletr\^{o}nica de Compras} (BEC), an electronic platform that processed over R\$7.4~billion for the State Health Secretariat and other agencies during our study period. The platform records every bid submitted, the identities of bidders and winners, the auction format, and the reference price set by the buyer. Two formats predominate: the \emph{convite} (sealed-bid tender) and the \emph{preg\~{a}o eletr\^{o}nico} (descending auction with real-time bidding). Our analysis focuses on preg\~{a}o auctions, which account for the majority of BEC spending and have the richest bid-level data.

The BEC catalog organizes items by group and class, totaling about 6{,}350 distinct product codes. Each unique procurement event is defined by a purchase-offer--item-code pair. This granularity allows us to track which firms bid on which products over time---a feature essential for identifying whether a firm enters new product markets after absorbing workers from a cartel competitor.

\subsection{CADE cartel convictions}

Brazil's federal competition authority, CADE, convicted firms in seven procurement cartels operating in S\~{a}o Paulo between 2009 and 2019, spanning pharmaceuticals, school meals, trash bags, solar water heaters, metropolitan rail, airport cafeterias, and school transport. Thirty-five firms were convicted, and their CNPJ roots can be matched to BEC records. Each case file specifies the colluding firms, the product market, and the conduct period, allowing us to identify precisely when cartel operations ended and the post-conduct labor reallocation began.

\subsection{RAIS: the employer--employee link}

The \emph{Rela\c{c}\~{a}o Anual de Informa\c{c}\~{o}es Sociais} (RAIS) is Brazil's compulsory annual census of formal employment. Every establishment must report each employment bond, identified by the worker's PIS number. We use the 2009--2017 vintages, restricting to establishments whose CNPJ root appears in the BEC firm universe. This allows us to track individual workers as they move from cartel firms to competitors after conduct periods end, and to characterize the movers by occupation, wage, tenure, and education.
